% Passes 5832--5839: normalizer/code/Pauli/all-field matrix closure.
\subsection*{Matrix-torsor normalizer, binary Reye code, and the dual-label Pauli quadric}

The affine matrix model of the selected $q=5$ Reye/Latin carrier admits a complete
normalizer calculation.  Put
\[
G=M_2(\mathbb F_2)_+\rtimes\bigl(GL_2(2)\times GL_2(2)\bigr),
\qquad |G|=576,
\]
acting on the $16$ matrices by $M\mapsto AMB^{-1}+X$.  Exhaustion of
$GL_4(2)$ gives
\[
\left|N_{GL_4(2)}(GL_2(2)\times GL_2(2))\right|=72,
\qquad
\left|C_{GL_4(2)}(GL_2(2)\times GL_2(2))\right|=1.
\]
The regular translation subgroup $T\cong C_2^4$ is $O_2(G)$, since
$G/T\cong S_3\times S_3$ has trivial $2$-core.  Thus $T$ is characteristic and
any permutation normalizing $G$ lies in $N_{S_{16}}(T)=AGL(4,2)$.  Consequently
\[
\boxed{|N_{S_{16}}(G)|=1152},
\qquad
\boxed{N_{S_{16}}(G)/G\cong C_2},
\]
with the nontrivial coset represented by matrix transpose.  This proves, in the
literal degree-$16$ carrier, that the factor-swap outer class is unique.

The characteristic-$2$ Reye kernel also becomes objectwise transparent.  For
$P=\{(w,x):0\ne w\in\mathbb F_2^2,\ x\in\mathbb F_2^2\}$ and
$L_M=\{(w,Mw):w\ne0\}$, the binary orthogonal kernel of the $12\times16$
point--line incidence is the two-weight code
\[
[12,4,6],\qquad W(z)=1+12z^6+3z^8.
\]
Its twelve minimum words are exactly the twelve heavy supports
\[
c_{\phi,\psi}(w,x)=\phi(x)+\psi(w),\qquad \phi\ne0,
\]
while the three weight-$8$ words are the $\phi=0$, $\psi\ne0$ fiber-pair
words.  Moreover this code is exactly the puncture of the binary simplex
$[15,4,8]$ code on $PG(3,2)$ obtained by deleting the projective line
$\{(0,x):x\ne0\}$.  This supplies a conceptual proof of the enumerator.  On the
saturated common-$W_9$ lattices the Smith form
\[
\operatorname{SNF}(R^T)=1^5\,2^2\,4^2
\]
has four even invariant factors and $2$-adic cokernel valuation six, so the
four-dimensional mod-$2$ defect is the integral lift of this same local Reye
kernel.  The heavy transform has the distinct Smith form
$1^2\,2^5\,4^2$ and must not be identified with it.

There is also now an object-level, rather than cardinality-only, bridge to the
published two-qubit projective-ring geometry.  Using the Pauli binary coordinates
$v=(x_1,z_1,x_2,z_2)$ for the $C_1,\ldots,C_{15}$ labeling of
Saniga--Planat--Pracna, their $9+6$ partition is the zero/one split of
\[
q_{\rm SPP}(v)=x_1z_1+x_2z_2+x_2.
\]
For a nonzero dual Fourier label
$Y=\bigl(\begin{smallmatrix}a&b\\c&d\end{smallmatrix}\bigr)$, the matrix-rank
split is the zero/one split of
\[
q_M(Y)=\det Y=ad+bc.
\]
There are exactly $72$ elements $L\in GL_4(2)$ satisfying
$q_{\rm SPP}(L(Y))=\det Y$ for every $Y$.  One explicit choice sends the
matrix-coordinate basis $(a,b,c,d)$ to the Pauli labels
\[
(XI,\ IZ,\ IY,\ ZI).
\]
It sends the nine nonzero rank-one labels exactly to the published Mermin
nine and the six units exactly to its six-point complement, while preserving
all $105$ polarized symplectic pairings.  The rank-one $3\times3$ grid becomes
\[
\begin{pmatrix}
C_{10}&C_{13}&C_7\\
C_{12}&C_{15}&C_9\\
C_{11}&C_{14}&C_8
\end{pmatrix},
\]
whose rows and columns are commuting triples.  This sharpens the prior boundary:
the two $9+6$ decompositions are isomorphic quadratic/symplectic geometries on
the \emph{dual Fourier labels}; it does not identify the $q=5$ cover points with
physical qubit states or observables.

The binary $1+9+6$ decomposition is the $r=2$ member of an all-field theorem.
For every finite field $\mathbb F_r$, let $R$ be the evaluation incidence
\[
R[(w,x),M]=1\iff x=Mw,
\qquad 0\ne w\in\mathbb F_r^2,\quad M\in M_2(\mathbb F_r).
\]
Additive Fourier labels split by matrix rank with multiplicities
\[
N_0=1,\qquad N_1=(r-1)(r+1)^2,\qquad
N_2=r(r-1)^2(r+1)=|GL_2(r)|.
\]
Then
\[
\boxed{\operatorname{im}R^T=\mathbf 1\oplus\mathcal F_{\mathrm{rank}\,1}},
\qquad
\boxed{\ker R=\mathcal F_{\mathrm{rank}\,2}},
\]
so
\[
\boxed{\operatorname{rank}R=1+(r-1)(r+1)^2}.
\]
The two nonzero squared singular values are $r^2(r^2-1)$ on the constant mode
and $r^2(r-1)$ on rank one.  The companion line--hyperplane disjointness relation
is
\[
D[M,(\phi,\psi)]=1\iff \psi=-\phi M.
\]
At $r=2$ every rank-one label has a unique nonzero factorization
$Y=\phi^T w^T$; for $r>2$ it has $r-1$ factorizations, producing an additional
point-side kernel of dimension $(r-2)(r-1)(r+1)^2$.  Exact Fourier replay at
$r=2,3,5,7$ agrees with these formulas; the proof extends to prime powers by
composing a nontrivial additive character with the field trace.

Two further finite consequences expose the same split from different directions.
First, the determinant phase is a self-dual bent quadratic:
\[
\boxed{\widehat{(-1)^{\det}}(Z)=4(-1)^{\det Z}}
\qquad (Z\in M_2(\mathbb F_2)).
\]
Second, the Cayley graph on additive $M_2(\mathbb F_2)$ with
$M\sim N$ iff $M-N$ is invertible is
\[
\operatorname{SRG}(16,6,2,2)=L_2(4).
\]
Its eight maximal $K_4$'s form two parallel classes, recovering intrinsic
$4\times4$ rook coordinates.  Its full automorphism group is
$S_4\wr C_2$ of order $1152$, and the resulting permutation group is exactly
the normalizer above.  Hence the previously unexplained order-$1152$ rook
symmetry comes from the six unit-difference directions, not from the Delsarte
graph of the Reye two-weight code.

\paragraph{Boundary.}
These are exact statements in finite affine geometry, binary coding, integral
Smith theory, Fourier analysis, graph theory, and two-qubit Pauli incidence.
The Pauli bridge lives on nonzero dual Fourier labels and does not supply a
$q=5$ physical-state embedding, dynamics, mass or coupling law, continuum limit,
or experimental prediction.  The local $[12,4,6]$ code is also not identified
with the global $q=5$ binary $[[156,26,6]]_2$ point-code CSS family merely because
their minimum distances coincide.
